\documentclass[10pt,a4paper]{amsart}
\usepackage[margin=19mm]{geometry}
\usepackage{amsmath,amssymb,mathtools}
\usepackage[scr=rsfs,cal=euler]{mathalfa}
\usepackage{xfrac}
\usepackage[unicode,hidelinks,bookmarks=false]{hyperref}
\newcommand{\ie}{i.e.\ }
\newcommand{\cf}{cf.\ }
\newcommand{\resp}{respectively\ }
\newcommand{\Subsection}{\subsection}
\providecommand{\nobreakdash}{\nobreak\hbox{-}}
\setlength{\emergencystretch}{2em}
\multlinegap0pt
\usepackage{mathtools}
\usepackage{amssymb}
\usepackage{xcolor}
\newtheorem{prop}{Proposition}
\newtheorem{rem}{Remark}
\newtheorem{cor}{Corollary}
\newtheorem{lem}{Lemma}
\newcommand{\G}{\mathbb{G}}
\renewcommand{\L}{\mathcal{L}}
\newcommand{\R}{\mathbb{R}}
\newcommand{\cg}{\Gamma}
\renewcommand{\l}{\lambda}
\DeclareMathOperator{\supp}{supp}
\title{Capacities Characterizing Removable Sets for Various Function Spaces in Carnot Groups}
\author{Zack Boone}
\date{Source: arXiv:2512.17167v1 (CC BY 4.0)}
\begin{document}
\maketitle

\noindent\emph{Source and licence.} Capacities Characterizing Removable Sets for Various Function Spaces in Carnot Groups. Version arXiv:2512.17167v1 (CC BY 4.0), distributed by arXiv under the Creative Commons Attribution 4.0 International licence, http://creativecommons.org/licenses/by/4.0/, as recorded in the arXiv licence metadata for this exact version and shown on the abstract page https://arxiv.org/abs/2512.17167v1. Full licence notice: \texttt{licenses/arxiv-2512.17167-CC-BY-4.0.txt}.\par\smallskip

\noindent\textit{Editorial scope.} Counted unit: the article's lem third\_func\_in\_lip (source0.tex lines 448--450). The article's complete original proof of it is delivered with it (source0.tex lines 451--471, one \texttt{proof} environment of 21 lines). No mathematics is written, repaired or re-derived: the statement and the proof are the article's own lines, in the article's own notation, and no retained line is substituted or re-flowed.  Retained uncounted as the source-local context the proof needs: lines 285--287; lines 289--291; lines 293--295; lines 371--373; lines 432--436; lines 581--583.  Nothing was omitted from any retained span, so the package carries no omission record.  No retained line contains a citation, so the wrapper carries no bibliography and no bibliography entry.  No retained line contains a figure environment, an image-inclusion command or drawing code, so no image is delivered and the package declares no asset dependency.  Editorial formatting only, in addition: an AMS \texttt{amsart} preamble that restates the article's own declarations and macros used by the retained lines, namely the environments \texttt{prop}, \texttt{rem}, \texttt{cor}, \texttt{lem} on separate counters, together with the standard packages those lines need.  The retained environments print in this document as Proposition 1, Remark 1, Corollary 1, Lemma 1 in order of appearance, whereas the article prints the same items with the numbering of its own full document.  The complete native source of the article as distributed in the arXiv e-print for this exact version is bundled unchanged as \texttt{sources/191295/source0.tex}.
\begin{prop}\label{hom-degree-of-deriv}
Let $f$ be a smooth homogeneous function of degree $m \in \R$ and $X$ be a linear differential operator of degree $n \in \R$. Then $Xf$ is a homogeneous function of degree $m-n$.
\end{prop}

\begin{rem}\label{deriv-of-fund-soln}
It is simple to check that the same proof which gives Proposition \ref{hom-degree-of-deriv} also works when we relax the smoothness condition to $f \in C^{\infty}(\G \backslash \{0\})$. In particular, if $k$ is a kernel of type $\l$ and $X$ a linear differential operator of degree $m$, then since $k$ is a smooth function on $\G \backslash \{0\}$, we have $X k$ is both smooth and homogeneous of degree $\l - Q - m$ on $\G \backslash \{0\}$. This follows since $k$ is smooth and homogeneous of degree $\l - Q$ on $\G \backslash \{0\}$.
\end{rem}

\begin{prop}\label{hom-func-bound}
Suppose $f$ is both homogeneous of degree $\l$ and continuous on $\G \backslash \{0\}$. Then, $$|f(p)| \leq C \| p\|^{\l}$$ for a constant $C>0$ which only depends on $f$ and $\G$. 
\end{prop}

\begin{cor}\label{bounded-deriv-is-lip}
Let $f : \G \rightarrow \mathbb{R}$ be a continuous function satisfying $\| \nabla_{\G} f\|_{L^{\infty}} \leq C$ for some finite constant $C$. Then $f$ is Lipschitz.  
\end{cor}

\begin{prop}\label{decay-of-one-deriv}
 Set $\xi(x) = f \L \varphi \ast X\cg(x)$ for $f \in \text{Lip}_{\text{loc}}$ and $\varphi \in C_0^{\infty}(B(a,\delta))$, and where $X$ is any element in $\mathcal{X}$. Then there exists $R < \infty$ such that \begin{align*}
     |\xi(x)|\leq 1 \quad \text{ for all }x\in B(a,R)^c.
 \end{align*}  
\end{prop}

\begin{lem}\label{second_func_in_lip}
Take $f\in \text{Lip}_{\text{loc}}$ and $\varphi \in C_0^{\infty}(B(a,\delta))$. Then the function $\Theta$ defined by $\Theta(x) := X_i(f X_i \varphi) \ast \cg(x)$ is continuous and satisfies $\| \nabla_{\G} \Theta \|_{\infty} < \infty$. In particular, Corollary \ref{bounded-deriv-is-lip} gives $\Theta \in \text{Lip}$.  
\end{lem}

\begin{lem}\label{third_func_in_lip}
Take $f \in \text{Lip}_{\text{loc}}$ and $\varphi \in C_0^{\infty}(B(a,\delta))$. Then the function $\Theta$ defined by $\Theta (x) := f\L \varphi \ast \cg(x)$ is continuous and satisfies $\| \nabla_{\G} \Theta \|_{L^{\infty}} < \infty$. In particular, Corollary \ref{bounded-deriv-is-lip} gives $\Theta \in \text{Lip}$.  
\end{lem}

\begin{proof}
%Like the proof of Lemma \ref{second_func_in_lip}, it suffices to show $\Theta \in C^1 \left( B(a,2\delta)\right)$. Take any $X \in \mathcal{X}$. Remark \ref{deriv-of-fund-soln} and Proposition \ref{hom-func-bound} say $|X \cg (p)|\lesssim \| p\|^{1-Q}$ for all $p \in \G \backslash \{0\}$. Let $g = f \L \varphi$. Then for any $x \in B(a,2\delta)$, left-invariance of $X$ yields, \begin{align}
%|X\Theta(x)| &\lesssim \|g\|_{\infty} \int_{B(a,\delta)} \| y^{-1}\cdot x\|^{1-Q}\,dy \nonumber \\ &\sim \sum_{j=0}^{\infty} \int_{ 2^{-j-1}(3\delta) < \| x^{-1}\cdot y\| \leq 2^{-j} (3\delta)} \|x^{-1}\cdot y\|^{1-Q}\,dy \nonumber \\ &\lesssim \sum_{j=0}^{\infty} (2^{-j-1}(3\delta))^{1-Q} \left| B(x, 2^{-j}(3\delta)) \right| \nonumber \\ &\sim \delta \sum_{j=0}^{\infty} 2^{-j} = \delta. \label{dct-bound-2}
%\end{align}
%A change of variables gives, \begin{equation}\label{dct-cont-2}
% X\Theta(x) = \int g(x \cdot p) X\cg(p)\,dp   
%\end{equation}
%and since $g$ is continuous, (\ref{dct-bound-2}) and (\ref{dct-cont-2}) allow the use of the dominated convergence theorem to conclude $\Theta \in C^1\left( \overline{B(a,\delta)} \right)$. 
Proposition \ref{decay-of-one-deriv} gives the existence of $R > 0$ such that \begin{align}
    |X \Theta(x)| \leq 1 \quad \text{ for all }x \in B(a,R)^c. \label{bound_outside_ball}
\end{align} 
Remark \ref{deriv-of-fund-soln} and Proposition \ref{hom-func-bound} say $|X \cg(p)|\lesssim \| p\|^{1-Q}$ for all $p \in \G \backslash \{0\}$
So for any $x \in B(a,R)$, integrating over annuli gives \begin{align*}
    |X \Theta(x)|  &\lesssim \| g\|_{L^{\infty}(B(a,\delta))} \int_{B(a,\delta)} \| y^{-1}\cdot x\|^{1-Q}\,dy  \\ &\leq \|g\|_{L^{\infty}(B(a, \delta))} \sum_{j=0}^{\infty} \int_{ 2^{-j-1}(R+\delta) \leq \| y^{-1} \cdot x\| < 2^{-j}(R+\delta)}\| y^{-1} \cdot x\|^{1-Q}\,dy \\ &\lesssim \| g\|_{L^{\infty}(B(a,\delta))}\sum_{j=0}^{\infty} (2^{-j} (R+\delta))^{1-Q} \left| B(x, 2^{-j}(R+\delta))\right| \\ &\sim \| g\|_{L^{\infty}(B(a,\delta))} (R+\delta)^{1-Q} \sum_{j=0}^{\infty} 2^{j(Q-1)} (R+\delta)^Q 2^{-jQ} \\ &= \| g\|_{L^{\infty}(B(a,\delta))}(R+\delta) \sum_{j=0}^{\infty} 2^{-j} \\ &\sim \| g\|_{L^{\infty}(B(a,\delta))}(R+\delta).
\end{align*}
Along with (\ref{bound_outside_ball}), this shows $\| \nabla_{\G} \Theta \|_{L^{\infty}} \lesssim 1$.

For continuity, simply note $\supp \L \varphi$ is compact, $\cg$ is in $L_{\text{loc}}^1$, and $f \L \varphi$ is a continuous function. So the dominated convergence theorem gives that $\Theta$ is continuous.


\end{proof}

\end{document}
